\documentclass[12pt, letterpaper]{article}
\usepackage[margin=1.2in]{geometry}
\usepackage{ebgaramond}
\usepackage[T1]{fontenc}
\usepackage{microtype}
\usepackage{setspace}
\usepackage{parskip}
\usepackage{titling}
\usepackage{fancyhdr}
\usepackage[utf8]{inputenc}

\setlength{\parindent}{2em}
\onehalfspacing

\pagestyle{fancy}
\fancyhf{}
\fancyhead[L]{\small\textit{Fifteen Minutes}}
\fancyhead[R]{\small\textit{NASA Mission Series v1.19}}
\fancyfoot[C]{\small\thepage}
\renewcommand{\headrulewidth}{0.4pt}

\title{\Huge\textbf{Fifteen Minutes}\\[0.5em]\large\textit{A short story from Mission 1.19 --- Mercury-Redstone 3}}
\author{NASA Mission Series\\[0.3em]\small Miles Tiberius Foxglove --- Mukilteo, WA}
\date{March 30, 2026}

\begin{document}
\maketitle
\thispagestyle{empty}

\vspace{1em}
\noindent\rule{\textwidth}{0.4pt}
\vspace{1em}

\noindent The thing about waiting is that it sharpens everything. Alan Shepard had been lying on his back in the Mercury capsule for four hours and fourteen minutes. The original hold had been for cloud cover. Then the inverter in the Redstone's power supply had shown an anomaly, and they'd pulled the gantry back to replace it. Then the cloud cover again. The count had been recycled twice. The medical sensors taped to his chest itched. His suit was pressurized. He needed to urinate and there was no provision for that in the flight plan because the flight plan assumed a man would not be lying on his back in a capsule for four hours before launch.

He keyed the mic. ``Why don't you fix your little problem and light this candle?''

The words became famous. What didn't become famous was the silence on the other end of the line --- the three seconds of dead air where the men in Mercury Control looked at each other and decided whether that was frustration or bravado or both. It was both. Shepard was a Navy test pilot. He had flown things that weren't supposed to fly and landed them on ships that were moving. He did not wait well.

\begin{center}
$\bullet\quad\bullet\quad\bullet$
\end{center}

At T minus zero, the Rocketdyne A-7 engine ignited. Seventy-eight thousand pounds of thrust from a single chamber burning liquid oxygen and ethanol --- essentially a refined V-2, the German rocket that had terrified London sixteen years earlier, now painted white and aimed at the sky with an American Navy commander sitting on top of it. The hold-down clamps released. The Redstone rose.

Shepard felt the acceleration build. One g. Two. Three. The vibration was rough but manageable --- he'd been in the centrifuge at Johnsville at twelve g, and this was a Sunday drive by comparison. The g-force climbed as the rocket consumed fuel and became lighter while thrust remained constant. Four g. Five. The sky outside the periscope was changing color, blue deepening toward indigo, indigo toward black. At six point three g, the Redstone's engine cut off.

Silence.

Not actual silence --- the instrument panel still hummed, the suit fans still whirred, the radio still crackled. But the roar was gone, and the vibration was gone, and the weight was gone. Shepard was floating. The straps held him in the couch but his arms drifted up, and the dust motes in the cabin hung motionless, and the pencil he'd tucked into the knee pocket of his suit slid free and hovered in front of his face.

Weightlessness. Zero g. The thing nobody could prepare you for because nobody had experienced it --- not in America, not until this moment. Gagarin had, twenty-three days earlier, but Gagarin was on the other side of an Iron Curtain and his debriefing notes were classified and all anyone in the West knew was that he'd gone up and come back and said it was beautiful.

\begin{center}
$\bullet\quad\bullet\quad\bullet$
\end{center}

The periscope activated. Shepard looked down and saw the Earth.

Later, in debriefing, he would describe it in the precise language of a test pilot: cloud patterns visible, Lake Okeechobee identifiable, the west coast of Florida clearly delineated, the Bahamas visible as light green patches against the dark blue Atlantic. He would note that the horizon was sharply defined, that the transition from Earth to space was a thin band of brilliant blue fading through darker blue to black. He would provide these observations in the same tone he'd use to describe the instrument panel readings, because that was the job: observe, report, fly the machine.

What he didn't say in debriefing --- what he said later, years later, when the weight of being first had settled into something he could carry without it changing his posture --- was that looking through the periscope at one hundred and eighty-seven point five kilometers above the Atlantic Ocean, he felt a pull in his chest that had nothing to do with g-forces. The Earth was there. All of it was there. Not a map, not a globe, not a photograph, but the actual thing, curved and blue and wrapped in atmosphere so thin it looked like it could be brushed away.

\emph{The Pacific giant salamander lives its entire life within a radius of twenty meters. Dicamptodon tenebrosus: the largest terrestrial salamander in North America, up to thirty centimeters from snout to tail tip, dark brown with pale marbling that looks like lichen on wet rock. It lives under logs in old-growth forest. Under one log. The same log, year after year. It hunts at night, in rain, within a body-length of its shelter. It breathes through its skin. Its world is twenty meters of damp forest floor, and within that world it is complete.}

\begin{center}
$\bullet\quad\bullet\quad\bullet$
\end{center}

He had five minutes. That was the arc of the thing --- a Redstone couldn't reach orbital velocity, so the trajectory was a parabola, up and over and back down, like a ball thrown from an impossibly strong arm. Five minutes and sixteen seconds of weightlessness at the top of that arc, and in those five minutes Shepard had a list of things to do.

First: manual attitude control. He grasped the hand controller --- a pistol-grip stick on the right armrest --- and pitched the capsule nose-down. The hydrogen peroxide thrusters fired in crisp, controlled bursts. The attitude indicator on the panel showed the capsule responding exactly as designed. He pitched back up. He yawed left, then right. He rolled. Each axis responded cleanly. The capsule went where he pointed it.

This was the point of the mission. Not the view. Not the weightlessness. Not the fact that an American had finally reached space twenty-three days after a Soviet had. The point was that a human being could operate a spacecraft. Could take the controls and fly the thing, make corrections, maintain orientation. The automatic system worked --- the ASCS had held the capsule steady during the powered ascent --- but Shepard was there to prove that the automatic system wasn't the only option. That when the machine's judgment wasn't enough, a pilot's judgment could take over.

He tested all three modes: full manual, fly-by-wire (where his inputs went through the electronic control system), and the automatic mode with manual override. Each worked. Each produced the expected response. The capsule was flyable.

\begin{center}
$\bullet\quad\bullet\quad\bullet$
\end{center}

S\o ren Kierkegaard wrote about the leap of faith. Not the version that got diluted into a greeting card sentiment about trusting the universe --- the actual thing, the thing Kierkegaard sweated over in Copenhagen in the 1840s. The paradox: you cannot know, before you leap, whether the landing will hold you. You can study. You can prepare. You can calculate the trajectory and check the math and test the equipment and run the simulations a thousand times. But the moment of the leap is the moment when preparation ends and commitment begins, and the gap between the two is where faith lives.

Shepard was not a philosopher. He was a Navy pilot from East Derry, New Hampshire, who had been selected for this job because he was very good at controlling machines in situations where the machines were trying to kill him. He didn't think about Kierkegaard. He thought about the checklist.

But the leap was real. The moment the hold-down clamps released and the Redstone started moving, every calculation became irrelevant. The trajectory was committed. The physics was in charge. All Shepard could do was ride the math and, when the time came, fly the machine. The gap between preparation and commitment --- the Kierkegaard gap --- was the three-inch space between Shepard's gloved hand and the attitude controller. He could study the controller. He could test it on the ground. But the moment he reached for it in weightlessness, one hundred and eighty-seven kilometers above the ocean, and it responded --- that was the leap. Not knowing. Doing.

\emph{The Pacific giant salamander does not deliberate. When the rain comes and the night comes and the forest floor is wet, it leaves the shelter of its log and hunts. It does not calculate the probability of finding prey or weigh the risk of predation. It hunts because it is a salamander and the rain is falling and the conditions are right. The leap, such as it is, is built into the animal. The commitment was made a hundred million years ago, when the first amphibian crawled onto land and decided to stay.}

\begin{center}
$\bullet\quad\bullet\quad\bullet$
\end{center}

Retrofire. Three solid rockets in the retrorocket package, mounted on the heat shield, firing in sequence. Each one pushed against his chest like a firm hand. The capsule decelerated, arc bending, trajectory steepening. The retro package jettisoned. The periscope retracted. The heat shield was now the leading edge, facing forward into the atmosphere the capsule was about to reenter.

The g-forces built. Not gradually, not gently. The capsule hit the atmosphere at the top of a steep ballistic arc --- steeper than any orbital reentry would be --- and the deceleration was savage. Six g. Eight. Ten. The edges of Shepard's vision grayed. Eleven g. His chest compressed under the load --- his body weighed nearly two thousand pounds, every heartbeat an act of force, the blood in his extremities draining toward his core. Eleven point six g. The peak. The hardest reentry any Mercury astronaut would ever experience.

He stayed conscious. The centrifuge training held. The gray edges stabilized, didn't close. His vision tunneled but didn't black out. He watched the instrument panel through a narrowing field and kept his hand on the controller because if the automatic system failed, if the attitude drifted, he was the backup. He was always the backup. That was what a pilot was for.

Outside the porthole: orange fire. The ablative heat shield was doing its job, fiberglass boiling away in a controlled burn that absorbed the kinetic energy of the suborbital arc. Three thousand degrees Fahrenheit at the shield surface, and on the other side of it, in the pressurized cabin, Shepard was sweating but alive, the temperature gauge reading a tolerable hundred and two degrees.

\begin{center}
$\bullet\quad\bullet\quad\bullet$
\end{center}

The drogue parachute deployed at twenty-one thousand feet. A jolt, a snap, the capsule swinging. Then the main parachute at ten thousand feet --- the ringsail canopy blooming above like a jellyfish, orange and white gores catching the sunlight, the capsule hanging beneath it in a slow pendulum swing. The Atlantic was blue below, very blue, and Shepard could see whitecaps.

Splashdown was hard. Not gentle, not a soft landing. The capsule hit the water like dropping a boulder off a bridge --- the impact shoved him down into the couch and then the capsule tipped sideways, rolled, righted itself. It bobbed. Water sloshed against the porthole. The recovery helicopter was already overhead; he could hear the rotors through the hull.

Fifteen minutes and twenty-eight seconds. From the pad to the ocean. One hundred and sixteen miles downrange. Peak altitude one hundred and sixteen point five statute miles. First American in space.

\emph{The Pacific giant salamander, after a night of hunting, returns to its log. The same log. It fits itself into the space beneath the bark --- the exact space, the depression its body has worn into the soil over years of returning --- and waits for the next rain. Twenty meters of world. Complete. The salamander does not wonder what lies beyond the twenty meters. It knows what lies within them, and that is enough. The log is old-growth Douglas fir, and the forest has been here since before the last ice age, and the salamander's ancestors have been hunting this same forest floor for fifty million years. The continuity is the point. The log endures. The salamander endures. The rain comes and the rain comes and the pattern holds.}

\begin{center}
$\bullet\quad\bullet\quad\bullet$
\end{center}

Kierkegaard also wrote about repetition. Not boredom --- not the deadening sameness of routine --- but the philosophical problem of whether a genuine experience can be repeated. Can you step into the same river twice? Can the thrill of the first time be recovered? Kierkegaard argued that repetition is a forward motion: not going back to what was, but finding again what was always there. The river has changed. You have changed. But the riverbed is still carved by water, and the current still pulls, and the act of stepping in is still the act of stepping in.

Shepard would fly again. Gemini was supposed to be his next mission, but an inner-ear condition grounded him. He fought for ten years to get back. He had surgery. He was reinstated. He commanded Apollo 14. He walked on the Moon. He hit a golf ball on the lunar surface --- a one-handed six-iron shot, in a pressure suit, at one-sixth gravity. It went, he estimated, miles and miles and miles.

But the fifteen minutes were first. The fifteen minutes were the leap. Everything after was repetition in Kierkegaard's sense: not a copy of the original, but a forward motion that carried the original within it. When Shepard stood on the Moon in 1971, the fifteen minutes of 1961 were still in his bones, still in the muscle memory of reaching for a hand controller in weightlessness and having the machine respond. The first time had carved the riverbed. Every subsequent step followed the current.

\emph{Dicamptodon tenebrosus. Tenebrosus: of the dark places. The species name is a description and a home address. Under the log, in the dark, in the wet forest of the Pacific Northwest where the canopy is so dense that even noon feels like twilight. The darkness is not absence. It is habitat. The salamander is built for it --- eyes adapted to low light, skin that breathes moisture, a metabolism tuned to cool, damp air. To pull the salamander from its darkness and put it in sunlight would not illuminate it. It would kill it. The darkness is where the animal lives. The darkness is home.}

Shepard came home. The helicopter lifted him from the capsule, the carrier USS \emph{Lake Champlain} was on station, the deck crew cheered. He walked across the flight deck in his silver pressure suit, visor up, grinning. He had been gone for fifteen minutes and twenty-eight seconds. He had been gone for four years of preparation. He had been gone for the three inches between his hand and the controller, the Kierkegaard gap, the leap and the landing and the discovery that the landing held.

He was home. The log endured. The salamander endured. The fifteen minutes were carved into the riverbed, and everything that followed --- Grissom's flight, Glenn's orbit, the Gemini program, Apollo, the Moon --- flowed along the channel they had cut.

\vspace{2em}
\begin{center}
\rule{3cm}{0.4pt}\\[1em]
\textit{Dedicated to S\o ren Kierkegaard (1813--1855)\\who wrote about leaps, and repetition, and what it means to commit}
\end{center}

\vspace{2em}
\begin{center}
\small
NASA Mission Series v1.19 --- Mercury-Redstone 3 --- May 5, 1961\\
Organism: \textit{Dicamptodon tenebrosus} (Pacific Giant Salamander) $\bullet$ Bird: Peregrine Falcon (degree 19)
\end{center}

\end{document}
